% ICLR 2027 pandoc template. Canonical copy lives here, in /latex-template at
% the repo root, alongside its style files (iclr2027_conference.sty/.bst/.bib,
% fancyhdr.sty, natbib.sty, math_commands.tex), sourced from
% https://media.iclr.cc/Conferences/ICLR2027/iclr-2027-style-files.zip
%
% The Obsidian "Enhancing Export" plugin's "ICLR 2027" preset is hard-coded to
% look up "neurips.tex" inside its own plugin folder's textemplate/ dir, so
% that directory holds symlinks back to the files here rather than copies --
% edit the files in /latex-template, not the symlinks.
% This template was tested with Pandoc 3.4. It should be backwards compatible with older version of pandoc..
\documentclass{article} % For LaTeX2e


% ready for submission (author names hidden, "Under review" header).
% Uncomment \iclrfinalcopy below for the camera-ready version.
\usepackage{iclr2027_conference,times}


\usepackage[utf8]{inputenc} % allow utf-8 input
\usepackage[T1]{fontenc}    % use 8-bit T1 fonts
\usepackage{hyperref}       % hyperlinks
\usepackage{url}            % simple URL typesetting
\usepackage{booktabs}       % professional-quality tables
\usepackage{amsfonts}       % blackboard math symbols
\usepackage{nicefrac}       % compact symbols for 1/2, etc.
\usepackage{microtype}      % microtypography
\usepackage{xcolor}         % colors
\usepackage{graphicx}
\usepackage{longtable} % Add support for Pandoc's longtable if needed
\usepackage{array}     % For table alignment improvements
\usepackage{amsmath}
\usepackage{textcomp}
\setlength{\LTcapwidth}{\textwidth} % To make captions fit within page width
\usepackage{calc} % arithmetic in table column-width specs (pandoc's multi-column
                   % longtable output needs this regardless of citeproc/csl-refs)

% pdfTeX (unlike xelatex/lualatex) can't set most Unicode math/arrow symbols
% directly even under utf8 inputenc + T1 fontenc -- map the ones the Obsidian
% source actually uses to their LaTeX math equivalents. Extend this list if a
% future draft introduces new Unicode symbols outside Latin-1.
\usepackage{newunicodechar}
\newunicodechar{≫}{\ensuremath{\gg}}
\newunicodechar{≪}{\ensuremath{\ll}}
\newunicodechar{→}{\ensuremath{\rightarrow}}
\newunicodechar{←}{\ensuremath{\leftarrow}}

\makeatletter
\newsavebox\pandoc@box
\newcommand*\pandocbounded[1]{% scales image to fit in text height/width
  \sbox\pandoc@box{#1}%
  \Gscale@div\@tempa{\textheight}{\dimexpr\ht\pandoc@box+\dp\pandoc@box\relax}%
  \Gscale@div\@tempb{\linewidth}{\wd\pandoc@box}%
  \ifdim\@tempb\p@<\@tempa\p@\let\@tempa\@tempb\fi% select the smaller of both
  \ifdim\@tempa\p@<\p@\scalebox{\@tempa}{\usebox\pandoc@box}%
  \else\usebox{\pandoc@box}%
  \fi%
}
\makeatother
\makeatletter
\def\maxwidth{\ifdim\Gin@nat@width>\linewidth\linewidth\else\Gin@nat@width\fi}
\def\maxheight{\ifdim\Gin@nat@height>\textheight\textheight\else\Gin@nat@height\fi}
\makeatother
% Scale images if necessary, so that they will not overflow the page
% margins by default, and it is still possible to overwrite the defaults
% using explicit options in \includegraphics[width, height, ...]{}
\setkeys{Gin}{width=\maxwidth,height=\maxheight,keepaspectratio}
% Set default figure placement to htbp
\makeatletter
\def\fps@figure{htbp}
\makeatother

\providecommand{\tightlist}{%
  \setlength{\itemsep}{0pt}\setlength{\parskip}{0pt}}
\title{Machine intelligence is idiosyncratic and incoherently
structured}


% Authors must not appear in the submitted version -- the iclr2027_conference
% style automatically prints "Anonymous authors / Paper under double-blind
% review" instead of the block below as long as \iclrfinalcopy stays commented
% out. Fill in real authors here for the camera-ready version.
\author{%
}


%\iclrfinalcopy % Uncomment for camera-ready version, but NOT for submission.
\begin{document}


\maketitle


\begin{abstract}
    We take a multi-order latent variable approach to intelligence in
    language models, similar to how psychometricians formulate
    psychological traits. Performance in every specific problem set is
    influenced by a domain-specific and a domain-agnostic latent factor.
    Using factor analysis as a dimension-reduction technique, we analyse
    13,251 published evaluation scores covering 1,618 language models
    across 456 different text-only benchmarks. Due to the super-sparse
    nature of the dataset, we triangulate our analysis across different
    data densifiers and imputation methods. A robust pattern across
    different modes of bias is that 1. factor patterns are only
    partially interpretable and often incoherent, 2. no silver-bullet
    ``general intelligence'' factor exists. Our findings goes against
    current endeavors of defining, identifying, and targeting general
    intelligence in language model development. It is not possible to
    develop a generally-intelligent language model by targeting single
    conceptual ability: general intelligence is only achievable by
    training on the first-order intelligence domains, but these are
    often partially idiosyncratic and not identifiable in practice.
  \end{abstract}


\hypertarget{introduction}{%
\section{Introduction}\label{introduction}}

The field of artificial intelligence is rife with theories and
conceptualization about what intelligence ``is'' (). With neural network
specifically, a longstanding issue is how the models are ``overfitting''
their training set, and thus their performance in various tasks
(``intelligence'', so to say) lack powerful generalizability that
cognitive entities like humans have. Thus, the development of advanced
intelligent systems rely on a working notion of intelligence as adaptive
and effective across diverse tasks: a singular notion of
``generalizability.''

There is, in this endeavor, an implicit and subtle assumption: it is
that in neural networks, much like humans, cognitive abilities are
structured according to a causal hierarchy. The word causal is important
here. A popular notion of general intelligence in the field is the
concept ``fluid intelligence'' borrowed from psychometric research,
defined loosely as the ability to perform well in any tasks regardless
of the content (i.e., a ``domain-free'' cognitive capacity). By
definition this construct holds a causal relation, where an increase of
fluid intelligence leads to an increase in the mastery of every other
cognitive task. In other words, domain-general ability \emph{precedes
and builds} domain-specific abilities.

It is, therefore, useful to ascertain whether it is plausible that
cognitive abilities in neural networks follow a coherent hierarchical
causal structure. An important implication is that if such structures
prove to be incoherent or idiosyncratic, then the entire program of
general intelligence as a silver bullet construct (e.g., ``AGI'') falls
apart and is entirely impossible: there is no efficient way to achieve
perfect-general intelligence without brute-forcing the training set to
cover the universe of all possible tasks, as there are no coherent
higher-order facets that researchers can target.

Why this causal hierarchy assumption is so rife in the first place is,
in our view, the application of both our theories and common sense about
our own intelligences to neural networks. No matter what one's stance on
the cognitive nature of AI is--whether or not they are conscious,
intelligent, or if they have internal representations proper--it is
reasonable to assert that AI are not humans. As such, the application of
top-down, human-based theories of intelligence to neural network models
are a category error. We cannot describe what machine intelligence is
like by overfitting theories of human intelligence on neural networks.

For this reason it is important to study the structure of machine
intelligence in a bottom-up manner. Borrowing from psychometrics yet
again, the application of exploratory factor analysis, a
dimension-reduction statistical method, allows us to summarize
performance in a series of abilities (measured through benchmarks) as
influenced by a higher-order latent variable in a plausibly causal
manner.

Prior work has tried applying factor analysis to benchmark scores. Ilica
\& Gignac (2024) have applied factor analysis to hundreds of models and
found that the models' performance fits a human-informed causal
structure. However, the study is flawed by their use of confirmatory
factor analysis, which is a hypothesis-testing procedure, and follows
the mistake of fitting a human-informed theory rather than understanding
model intelligence from the bottom-up. It is also severely flawed in
that a large proportion of their benchmarks are simply variants of the
popular MMLU benchmark (), heightens the risk of common-method bias
polluting the model fit.

Another relevant work is Burnell et al.~(2023), which correctly applied
EFA to model benchmark scores. They found that LLM abilities are
hierarchically structured under three major factors. However, they did
not do higher-order factor analysis, where EFA is run on the resulting
factor loadings, thus allowing us to understand to what extent are
latent factors influenced by a presumable g-factor.

A common weakness of both papers, though, are the relatively few number
of benchmarks or variables derived thereof. Ilica \& Gignac (2024) fit a
model with 20 variables, while Burnell et al.~(2023) uses 23 variables.
Most of these benchmarks are popular benchmarks measuring ``smartness''
for a lack of better term, like mathematics, academic knowledge, and
common-sense reasoning.

Within the context of prior work, our study presents an unprecedentedly
large factor analysis of benchmark scores. The curated archive holds
2,014 models and 624 benchmarks; the text-only subset all analyses run
on comprises 1,618 models and 456 benchmarks over 13,251 evaluation
scores. Crucially, our pool of benchmarks covers a highly diverse set of
tasks, including those quite outside of the mainstream. Covering only
few popular benchmarks like prior work does is problematic. First, given
their popularity, these benchmarks may well be correlated due to what we
call a common-investment bias: there is a high chance their correlations
substantially polluted by the fact that organizations expend more
efforts to perform well in said benchmarks. Second, there should be no
discrimination as to what tasks are important and which aren't: whether
a task appears miscellaneous is also no reason to exclude them. A useful
analogy is that reaction time in humans is correlated to intelligence
(Kranzler \& Jensen, 1989): no matter how seemingly unimportant a task
may be, it may provide useful information that is entirely unintuitive
to our subjective understanding.

\hypertarget{methodology}{%
\section{Methodology}\label{methodology}}

We ask whether the covariance among language models' published benchmark
scores admits a low-dimensional latent structure, and in particular
whether it supports a global general factor \(G\) in the sense of
\url{Appendix\#Definitions}. The empirical object is a \textbf{model ×
benchmark score matrix} assembled from public evaluation records. That
matrix is extremely sparse and its missingness is not at random:
widely-known models are evaluated on widely-used benchmarks, while
obscure benchmarks co-occur with almost nothing. Our methodology is
therefore organised around a single principle --- \textbf{no single
repair of the matrix is trustworthy on its own} --- so we run a
cross-product of deliberately different repairs and treat their
agreement or disagreement as the finding, rather than selecting one
recipe and reporting its output as the answer.

The pipeline has five stages: \emph{collection} → \emph{scoping} →
\emph{aggregation} → \emph{densification} → \emph{completion} →
\emph{factor analysis}.

All analyses reported here use the \textbf{text-only} subset of the
corpus (see \protect\hyperlink{Modalityux20scope}{\#Modality scope});
the multimodal-inclusive corpus is used only as a contrast condition
where explicitly stated.

\hypertarget{dataset}{%
\paragraph{Dataset}\label{dataset}}

\hypertarget{sources}{%
\subparagraph{Sources}\label{sources}}

Scores are transcribed from public evaluation records, not re-run by us.
We draw from four kinds of source, in descending order of volume: (i)
large curated evaluation suites, (ii) aggregate leaderboards, (iii)
benchmark-specific leaderboards, and (iv) primary papers. Concretely,
the corpus draws on the \textbf{Stanford HELM} family (Classic, Lite,
Safety, Reasoning, MedHELM, SEA-HELM, Arabic, ThaiExam, EWoK, TORR,
Finance), the \textbf{HuggingFace Open LLM Leaderboard} (v1 and v2),
\textbf{Papers With Code} evaluation tables, \textbf{Kaggle AI
Benchmarks}, \textbf{Chatbot Arena / LMArena}, \textbf{llm-stats.com},
\textbf{Artificial Analysis}, \textbf{Vellum}, and \textbf{LiveBench},
together with benchmark-specific leaderboards (e.g.~BigCodeBench,
CRUXEval, SWE-bench, BFCL/Gorilla, VMLU, SEA-LION) and primary arXiv/ACL
papers reporting original evaluations. Table 1 gives the composition of
the text-only corpus by source family;
\href{Appendix-Methods\#A\%20Data\%20sources\%20and\%20extraction}{Appendix A}
lists every named source and the extraction route used for each.

\textbf{Table 1.} Composition of the text-only corpus by source family
(13,251 result rows over 456 benchmarks and 1,618 models). Families are
aggregated from the recorded \texttt{source\_organization} field.

\begin{longtable}[]{@{}lrr@{}}
\toprule\noalign{}
Source family & Result rows & Distinct benchmarks \\
\midrule\noalign{}
\endhead
\bottomrule\noalign{}
\endlastfoot
Stanford HELM (11 sub-leaderboards) & 4,942 & 138 \\
HF Open LLM Leaderboard (v1 + v2) & 4,529 & 12 \\
Papers With Code & 1,378 & 151 \\
Kaggle AI Benchmarks & 844 & 26 \\
Primary papers (arXiv / ACL / journals) & 430 & 69 \\
Other named leaderboards and papers (26 sources) & 300 & 37 \\
Unattributed & 295 & 40 \\
Chatbot Arena / LMArena & 202 & 1 \\
llm-stats.com & 121 & 11 \\
Vellum & 84 & 7 \\
Artificial Analysis & 71 & 2 \\
LiveBench & 55 & 1 \\
\end{longtable}

\hypertarget{sources-partition-the-matrix}{%
\subparagraph{Sources partition the
matrix}\label{sources-partition-the-matrix}}

One consequence of this composition deserves stating before any of the
machinery, because it constrains what the analysis can recover.
Benchmark columns are largely \emph{owned} by a single evaluation suite,
and the suites evaluate near-disjoint sets of models. Measuring
co-observation between columns, grouped by the source that supplies most
of each column's rows:

\begin{longtable}[]{@{}
  >{\raggedright\arraybackslash}p{(\columnwidth - 4\tabcolsep) * \real{0.2727}}
  >{\raggedleft\arraybackslash}p{(\columnwidth - 4\tabcolsep) * \real{0.3636}}
  >{\raggedleft\arraybackslash}p{(\columnwidth - 4\tabcolsep) * \real{0.3636}}@{}}
\toprule\noalign{}
\begin{minipage}[b]{\linewidth}\raggedright
column pair
\end{minipage} & \begin{minipage}[b]{\linewidth}\raggedleft
median models observed on both
\end{minipage} & \begin{minipage}[b]{\linewidth}\raggedleft
pairs with \textless2 shared models (correlation not estimable)
\end{minipage} \\
\midrule\noalign{}
\endhead
\bottomrule\noalign{}
\endlastfoot
HELM × HELM & 4 & 33 \% \\
Open LLM Leaderboard × Open LLM Leaderboard & 160 & 0 \% \\
\textbf{HELM × Open LLM Leaderboard} & \textbf{0} & \textbf{87 \%} \\
\end{longtable}

The matrix is therefore closer to two weakly-connected blocks than to
one sparse whole: within the Open LLM Leaderboard block, correlations
rest on \textasciitilde160 shared models; across the two blocks, 87 \%
of column pairs cannot support a correlation estimate at all. This is
not a curation artefact we can repair --- it is a direct consequence of
which models each leaderboard chooses to evaluate --- and it means any
general factor recovered here is at risk of being a within-block factor.
We return to it when interpreting the results, and it is the structural
reason the cross-densifier comparison matters: the column-primary peel
retains famous benchmarks, which are disproportionately the
densely-connected block.

Two further properties matter downstream. First, source volume is highly
unequal: two source families supply two-thirds of all records, so the
\emph{missingness pattern} is largely inherited from those two suites'
model selection policies. Second, benchmark breadth and row volume are
anti-correlated --- Papers With Code contributes the most distinct
benchmarks but relatively few rows per benchmark, while the Open LLM
Leaderboard contributes 12 benchmarks with very deep model coverage.
This is precisely the structure that makes the densification choice in
\protect\hyperlink{Sparsityux20Handling}{\#Sparsity Handling}
consequential rather than cosmetic, and it is what motivates the
per-benchmark influence analysis in
\url{Analysis\#Robustness\%20and\%20sensitivity\%20analyses}.

\hypertarget{collection-protocol}{%
\subparagraph{Collection protocol}\label{collection-protocol}}

Collection follows a \textbf{strict source-verification} rule: a field
is populated only if the benchmark's authors or the evaluator explicitly
documented it. No value is inferred from a plausible default --- we
never assume, for example, that an open-weights model was evaluated with
a particular inference stack, or that an undocumented decoding
configuration was greedy. Undocumented fields are left blank. The one
class of exception is a small set of deductive rules that cannot be
wrong given the record itself (e.g.~a closed model accessed over a
vendor API cannot have been run locally); these are enumerated in
\href{Appendix-Methods\#A\%20Data\%20sources\%20and\%20extraction}{Appendix A}.

Records are stored in three linked tables --- \texttt{benchmarks} (one
row per benchmark), \texttt{models} (one row per model), and
\texttt{results} (one row per model × benchmark × evaluation setup).
Both \texttt{benchmarks} and \texttt{models} additionally carry a
\texttt{release\_date} (year-month) and a \texttt{release\_date\_source}
recording \emph{how} that date was established. The second column is not
bookkeeping: the sources differ sharply enough in reliability that
treating the date field as uniform would be a category error, and the
tiers are ordered so that an analysis can filter on them
(\href{Appendix-Methods\#A.5\%20Release-date\%20provenance}{Appendix A.5}).

The governing principle is that \textbf{a dating source is trustworthy
only when the model's identity is \emph{given} rather than inferred}. An
arXiv identifier encodes its submission month exactly; a HuggingFace
repository whose name \emph{is} the model reports its own creation
timestamp; a publisher's launch announcement names what it is
announcing. By contrast, any procedure that must normalise a model name
before matching it discards precisely the tokens that separate a model
from its family, and so systematically mis-dates a version to its
family's launch. We adopted this rule after two search-based sources
failed it on measurement rather than on principle, and after finding the
same error already present in the corpus: twenty models sat at GPT-4's
launch month, among them \texttt{GPT-4.1\ mini} and
\texttt{GPT-4.1\ nano}, which postdate it by two years.

Coverage is 2,007/2,014 models (99.7 \%) and 623/624 benchmarks (99.8
\%). The two tables are \textbf{not} of comparable quality, and the
near-identical coverage figures conceal this: 78 \% of model dates now
rest on an exact identifier, a repository timestamp, or a verified
announcement, against 0.5 \% of benchmark dates, 85 \% of which remain
on the weakest tiers. Model dates are at month precision in 97.3 \% of
rows; benchmark dates in 73.4 \%. Any temporal analysis should therefore
be run on the model axis and treat the benchmark axis as provisional
(\href{Appendix-Methods\#L\%20Known\%20limitations\%20and\%20deviations}{Appendix L.8}).
No stage of the pipeline reads \texttt{release\_date} --- densification,
completion and factoring operate on the score matrix alone --- so none
of the results reported here depend on it.

Multiple scores for the same model--benchmark pair are \textbf{kept as
separate rows} whenever they differ in evaluation setup, evaluator, or
language, rather than being averaged at collection time; they are
reconciled only at the aggregation step
(\protect\hyperlink{Aggregation}{\#Aggregation}), where the choice is
explicit and reversible. Scores are normalised to a 0--100 scale, with
documented exceptions for metrics that have no natural percentage
interpretation (perplexity, bits-per-byte, Elo); metric direction is
recorded rather than used to rescale
(\href{Appendix-Methods\#C\%20Normalisation\%20rules}{Appendix C}).

Inclusion criteria are applied at both axes. A \textbf{model} is
included if it is a generative language model that accepts arbitrary
prompts; encoder-only classifiers, narrow task-specific systems
(dedicated MT/ASR/TTS models), and undocumented community uploads are
excluded, and different \emph{setups} of one model (context length,
reasoning effort, prompting scheme) are recorded as setup attributes
rather than as distinct models. A \textbf{benchmark} is included if it
has at least one in-scope result row; zero-result stubs are removed. We
deliberately impose no relevance filter on benchmark content: a task is
not excluded for appearing miscellaneous or unrelated to
``intelligence'', for the reason given in the introduction. Full
criteria, and the review procedure applied to every borderline case, are
in
\href{Appendix-Methods\#B\%20Inclusion\%20and\%20exclusion\%20criteria}{Appendix B}.

Every edit to the corpus is followed by an automated integrity pass:
referential integrity in both directions, absence of orphan rows on
either axis, duplicate detection on a composite evaluation-identity key,
and a flag for benchmarks with implausibly few rows. Duplicate detection
distinguishes \emph{pure redundancy} (the same evaluation transcribed
twice) from \emph{conflicts} (different scores for what should be the
same evaluation); conflicts are resolved by source-trust tier and
recency, never silently averaged
(\href{Appendix-Methods\#C\%20Normalisation\%20rules}{Appendix C}).

\hypertarget{modality-scope}{%
\subparagraph{Modality scope}\label{modality-scope}}

All analyses in this paper are performed on a \textbf{text-only} subset.
The motivation is construct validity: a vision- or audio-conditioned
benchmark scores a model's perceptual front-end at least as much as its
language ability, so including such benchmarks mixes two different
latent spaces into one covariance matrix. A pilot comparison
additionally showed that removing them changes held-out predictive
accuracy and factor counts materially, i.e.~the distinction is not
cosmetic.

The restriction is applied by an auditable classifier rather than by
hand. Each benchmark's identifier and free-text metadata (category,
subcategory, task type, domain, name, description) are pooled and
matched against a fixed vocabulary of non-text modality terms under
\textbf{whole-word} matching, which avoids the substring false positives
(e.g.~\emph{vision} inside \emph{revisionism}) that a naive
contains-check produces. Because leaderboard metadata is itself
unreliable, the classifier is bracketed by two manually curated,
mutually disjoint override sets consulted before the pattern match: an
\textbf{allow-list} of benchmarks whose metadata suggests a non-text
modality but which are text-only on inspection (e.g.~a music benchmark
in ABC notation, a clinical-note task whose ``spoken dialogue'' is
supplied as a transcript), and a \textbf{deny-list} of benchmarks
confirmed non-text despite absent or mislabelled metadata (e.g.~an entry
named ``Chinese Multilingual MMLU'' whose rows in fact cite CMMMU, a
multimodal benchmark, and were scored on vision-language models). Every
override carries a written justification and the evidence used
(\href{Appendix-Methods\#D\%20Text-only\%20classifier}{Appendix D}).

Classification is followed by a \textbf{cascade removal}: each non-text
benchmark is dropped together with all of its result rows, and any model
left with zero remaining results is dropped in turn; the integrity
checks are re-run on the result. The derived copy is regenerated from
the canonical tables by script and is never hand-edited, so re-running
it after any corpus change or any revision to the pattern and override
sets is a single deterministic operation. This step removes 121 of 624
benchmarks (2,100 result rows) and, by cascade, 345 models, leaving 503
benchmarks and 1,669 models.

\hypertarget{score-redundant-benchmark-splits}{%
\subparagraph{Score-redundant benchmark
splits}\label{score-redundant-benchmark-splits}}

Public leaderboards frequently publish one benchmark as several
near-identical columns --- version-dated re-releases, difficulty or
subset variants, per-language splits of a translated test set, or the
same benchmark re-imported from a second source. Each such column enters
the matrix as a nominally distinct benchmark, and a factor analysis will
duly recover a ``factor'' that is nothing more than one benchmark's
identity replicated \(k\) times. Because our factor-count and
general-factor estimates are exactly the quantities such duplication
inflates --- and because this is one of the criticisms we level at prior
work --- we prune these before analysis.

Pruning is evidence-driven, not name-driven. For each suspected family
we compute the full pairwise Pearson correlation among its columns over
the models evaluated on both, and decide per family: a family is
collapsed to a single representative only when its columns are
near-perfectly correlated across a non-trivial shared model set. Where
correlations reveal genuine structure instead of redundancy, the family
is kept intact --- and in two of the eight audited families it was kept,
wholly or partly, for that reason. This yields seven removals in total
(version splits, dialect splits, difficulty and shot variants,
per-language splits of a translated benchmark, a cross-source re-import,
a corrupted composite identifier, and one partially redundant
national-exam trio), for 47 benchmark identifiers and 2,216 result rows.
Every decision, with its correlation statistics and the reasoning for
keeping or dropping, is recorded in
\href{Appendix-Methods\#E\%20Score-redundancy\%20pruning}{Appendix E}.

A separate, earlier pass removed \emph{translation duplicates} at corpus
level --- benchmarks that are literal translations of an original
already present --- while retaining multilingual benchmarks whose
per-language content is independently sourced. That distinction is a
content judgement made against each benchmark's source paper, not a
heuristic
(\href{Appendix-Methods\#C\%20Normalisation\%20rules}{Appendix C}).

After both passes, and after the canonical-metric filter and scale fix
described below remove a further 1,463 rows, the text-only corpus
contains \textbf{456 benchmarks, 1,618 models, and 13,251 result rows}
--- from a canonical archive of 624 benchmarks, 2,014 models and 19,030
result rows.

\hypertarget{aggregation}{%
\subparagraph{Aggregation}\label{aggregation}}

Factor analysis requires one score per model--benchmark cell, so the
multiple evaluation rows retained at collection time are averaged within
each (model, benchmark) pair. Averaging is deliberately placed here,
after all identity cleanup, so that it never masks a duplicate that
should have been removed.

Model identity is then resolved at \textbf{two granularities}, run as
parallel conditions throughout the rest of the pipeline:

\begin{itemize}
\tightlist
\item
  \textbf{\texttt{all\_standard}} --- variant-level. Source-specific
  model identifiers are normalised (organisation prefixes stripped,
  release dates and checkpoint stamps removed, context-length and
  reasoning-effort tags dropped, parameter counts canonicalised) so that
  different spellings of the same released variant collapse together,
  while genuinely different variants (sizes, generations, named tiers)
  stay distinct.
\item
  \textbf{\texttt{all\_aggressive}} --- family-level. Every model is
  collapsed to its base family token, so all sizes and generations of a
  family form one row.
\end{itemize}

The two strategies trade sample size against row homogeneity: the
standard collapse preserves more rows but leaves each row thinly
observed; the aggressive collapse produces far fewer, much
better-observed rows at the cost of treating a 7B and a 405B model of
one family as one entity. Neither is correct \emph{a priori}, which is
why both are carried forward. The token-level rules are given in
\href{Appendix-Methods\#F\%20Model-identity\%20collapse}{Appendix F}.

\hypertarget{metric-selection}{%
\subparagraph{Metric selection}\label{metric-selection}}

A benchmark is frequently reported under several metrics --- 92 of ours
were --- and the aggregation above would average them into one cell.
That is not merely untidy: on \texttt{truthfulqa} the cell would mix a
``\% informative'' rate with a BLEU \emph{difference} score, which can
be negative. Worse, \textbf{which} metric a model received is largely
determined by which leaderboard evaluated it, so the resulting column
carries variance attributable to its source rather than to capability.
On \texttt{truthfulqa}, 325 models are scored by accuracy (mean 46.9)
and 67 by exact match (mean 27.5), with the two populations essentially
disjoint; on \texttt{ifeval} the corresponding gap is 46 points. Because
no model is scored under both, the offset cannot even be estimated and
removed from within the data.

We therefore keep exactly \textbf{one metric per benchmark}. Metric
names are first normalised for case and whitespace --- without which
\texttt{Accuracy} and \texttt{accuracy} count as rivals and
\texttt{gsm8k} would forfeit 149 rows to a capitalisation difference ---
then a small curated alias map merges verified spelling variants of one
measurement (e.g.~\texttt{bits\ per\ byte} and \texttt{bpb}). Where a
genuine choice remains, we keep the metric covering the most distinct
models, so the widest comparable population survives; a per-benchmark
override list handles cases where coverage alone chooses badly. This
drops 1,420 result rows and 706 model-cells.

Two residual defects sit \emph{below} the metric name and need separate
treatment. \texttt{gpqa} carried one metric name covering two
incompatible conventions: 447 of its 454 rows are Open LLM Leaderboard
v2 \textbf{normalised} accuracy, in which the random-chance baseline is
mapped to zero (range 0--24.9), while the remaining 7 are raw accuracy
from papers (range 39--94). Since correlations are unchanged by a linear
rescaling of an entire column, a normalised column is perfectly usable
\emph{provided every row shares the convention}; we therefore keep the
447 and drop the 7, rather than back-transforming on an assumed formula.
The caveat is that 13 \% of those rows sit exactly at the clamp, so the
column under-discriminates among weak models. \texttt{elephant} is
removed outright: its metric field holds model configurations rather
than metrics, so no choice among them measures anything.
(\texttt{vectara}, whose two metrics were complements, and
\texttt{pwc\_lambada}, which mixed accuracy with perplexity, are both
resolved by the metric filter itself.)

Finally, benchmarks observed for only one model are dropped, since a
single observation contributes no covariance. The resulting matrices
are:

\textbf{Table 2.} Aggregated model × benchmark matrices, text-only
corpus. \emph{Provisional --- these matrices predate the
score-redundancy pruning and the canonical-metric filter, and are
superseded by the recomputed values recorded above; they are retained
only until the imputation and factoring results are re-run against the
current corpus, so that Tables 2, 3 and the Results tables can be
replaced together
(\href{Appendix-Methods\#L\%20Known\%20limitations\%20and\%20deviations}{Appendix L.3a}).}

\begin{longtable}[]{@{}lrrrr@{}}
\toprule\noalign{}
Collapse strategy & Models & Benchmarks & Observed cells & Density \\
\midrule\noalign{}
\endhead
\bottomrule\noalign{}
\endlastfoot
\texttt{all\_standard} (variant-level) & 1,310 & 455 & 13,888 & 2.33
\% \\
\texttt{all\_aggressive} (family-level) & 350 & 431 & 5,358 & 3.55 \% \\
\end{longtable}

\hypertarget{sparsity-handling}{%
\paragraph{Sparsity Handling}\label{sparsity-handling}}

At 2--4 \% observed, neither matrix admits a classical factor analysis:
the correlation matrix cannot be estimated by listwise deletion, since
there are no complete cases, and pairwise-complete estimation leaves
many benchmark pairs with zero or near-zero co-observation. The
missingness is moreover \textbf{not at random} --- a benchmark is
missing for a model precisely because that model was not considered
interesting enough to evaluate on it, which is itself a function of the
latent ability we are trying to measure.

\hypertarget{densification}{%
\subparagraph{Densification}\label{densification}}

We therefore construct \textbf{densified sub-matrices}, and we construct
several of them with deliberately opposed biases rather than one
``best'' one. All three densifiers greedily peel the matrix toward a
common target density (10 \%), differing only in which axis they
sacrifice:

\begin{itemize}
\tightlist
\item
  \textbf{C (column-primary)} --- repeatedly drop the least-observed
  \emph{benchmark}, then any model left empty. Retains famous benchmarks
  with wide model coverage.
\item
  \textbf{R (row-primary)} --- repeatedly drop the least-observed
  \emph{model}, then any benchmark left empty. Retains a broad benchmark
  set, including obscure ones, over a small set of heavily-evaluated
  models.
\item
  \textbf{S (symmetric)} --- at each step drop whichever marginal has
  the lowest \emph{fill-rate}, privileging neither axis.
\end{itemize}

After peeling, a floor is enforced on both axes so that every retained
model and benchmark has at least two observed scores; columns with zero
variance among observed values are also dropped, since they carry no
correlational signal. The peel targets density only. We deliberately do
\textbf{not} optimise pairwise overlap or positive-definiteness, because
those are the success criteria of particular downstream estimators ---
optimising them here would tilt the comparison in
\protect\hyperlink{Matrixux20completion}{\#Matrix completion} toward the
methods that need them. The algorithm is given in
\href{Appendix-Methods\#G\%20Densification\%20algorithm}{Appendix G}.

\textbf{Table 3.} Densified matrices (text-only corpus). ``Retained'' is
the fraction of observed cells surviving the peel.

\begin{longtable}[]{@{}lllrr@{}}
\toprule\noalign{}
Densifier & Strategy & Shape & Density & Retained \\
\midrule\noalign{}
\endhead
\bottomrule\noalign{}
\endlastfoot
C & \texttt{all\_standard} & 735 × 90 & 13.08 \% & 62.3 \% \\
C & \texttt{all\_aggressive} & 226 × 115 & 12.58 \% & 61.0 \% \\
S & \texttt{all\_standard} & 652 × 164 & 10.01 \% & 77.1 \% \\
S & \texttt{all\_aggressive} & 131 × 344 & 10.01 \% & 84.2 \% \\
R & \texttt{all\_standard} & 227 × 382 & 10.85 \% & 67.8 \% \\
R & \texttt{all\_aggressive} & 106 × 404 & 10.50 \% & 83.9 \% \\
\end{longtable}

The three densifiers span the aspect-ratio space: C yields tall matrices
(models ≫ benchmarks), R yields wide ones (benchmarks ≫ models), and S
sits between. Since the identifiability of a factor solution depends
strongly on that ratio, disagreement between C, S, and R is itself a
diagnostic, and we report it as one. The undensified matrix is
additionally carried through the pipeline as a \texttt{raw} contrast
level, analysed without imputation.

\hypertarget{matrix-completion}{%
\subparagraph{Matrix completion}\label{matrix-completion}}

Every densified matrix is still \textasciitilde90 \% missing, so a
completion step is required before factoring. We use \textbf{three
families} of method with different --- in places incompatible ---
assumptions, so that a structure recovered by all of them is unlikely to
be an artefact of any one.

\textbf{Cell-level imputation} estimates the missing entries directly:
\textbf{SoftImpute} (nuclear-norm-penalised low-rank completion by
iterative soft-thresholded SVD; assumes a low-rank signal plus noise,
and is our primary cell-level method); \textbf{k-NN} (each missing cell
filled from the \(k\) most similar models --- an assumption-light
baseline with no low-rank, linearity, or normality assumption);
\textbf{missForest} (iterative random-forest imputation, nonparametric,
able to capture nonlinear dependence the low-rank methods cannot
represent); and \textbf{MICE} (multiple imputation by chained equations,
where factoring receives the mean of the \(m\) completions while the
spread across completions is the only natively probabilistic uncertainty
estimate available to us).

\textbf{Correlation-level recovery} exploits the fact that factor
analysis needs a correlation matrix, not a data matrix. When
observations are too sparse to complete cells reliably, the correlation
structure may still be recoverable --- a substantially weaker
requirement. This family estimates the benchmark × benchmark correlation
matrix directly and never claims to know individual cells:
\textbf{OneSidedMC} (Cao, Liang \& Valiant, 2023) recovers the
benchmark-space right singular vectors from pairwise products of
co-observed scores, yielding an estimate \(\hat{\Theta}\) of the
benchmark covariance; \textbf{SoftImpute-corr}, \textbf{OptSpace}
(Keshavan, Montanari \& Oh, 2010), and \textbf{USVT} (Chatterjee, 2015)
apply matrix-completion estimators to the \emph{observed pairwise
correlation matrix}, whose missing entries are exactly the benchmark
pairs that were never co-observed; and two structured completions target
positive-definiteness directly --- a \textbf{maximum-determinant} SDP
completion, which maximises \(\log\det\Sigma\) subject to
\(\Sigma \succeq 0\) and to each observed correlation lying within a
per-pair Fisher-\emph{z} confidence band scaled to that pair's
co-observation count, and a \textbf{Gaussian graphical model} MLE
completion over the observed-pair graph.

Because these methods produce a correlation matrix rather than data, two
shared devices make them commensurable with the cell-level family.
First, the completed correlation matrix is symmetrised and projected to
the nearest valid (positive definite, unit-diagonal) correlation matrix,
so every principal submatrix is invertible. Second, a
\textbf{covariance-matched surrogate} data matrix is synthesised whose
sample covariance equals the recovered correlation matrix, on the
original column scale, and that surrogate is handed to the identical
factoring code used for every other method. The surrogate is explicitly
\emph{not} an estimate of the real cells; it is a device for passing a
covariance structure through a data-matrix interface, and no per-model
quantity computed from it is interpretable
(\href{Appendix-Methods\#H\%20Completion\%20methods}{Appendix H}).

\textbf{No-imputation baselines} bound how much of any recovered
structure is manufactured by imputation. The undensified matrix is also
factored with no completion at all, from a pairwise-complete correlation
matrix. Because that matrix has undefined entries (never co-observed
pairs) and is generally indefinite, two treatments are compared: filling
with the mean off-diagonal correlation, and filling with zero (treating
absent co-observation as absent association). Both are then PSD-smoothed
before factoring
(\href{Appendix-Methods\#H\%20Completion\%20methods}{Appendix H}).

\hypertarget{evaluating-the-completion}{%
\subparagraph{Evaluating the
completion}\label{evaluating-the-completion}}

All methods, in all three families, are scored on \textbf{one held-out
metric}, which is what makes them comparable at all. We mask
\textasciitilde20 \% of the observed cells using a
\textbf{column-stratified} split --- sampling within each benchmark, so
that no benchmark is absent from the evaluation set and high-frequency
benchmarks cannot monopolise it --- subject to leaving at least two
training observations in every column. Columns are standardised using
\textbf{training-cell moments only}; rows are never standardised,
because rows are models and row-standardisation would remove exactly the
between-model level differences that a general factor consists of.

Held-out cells are scored in standard-deviation units against a baseline
that predicts each column's \emph{training} mean:

\[\text{RMSE} = \sqrt{\overline{(\hat z - z)^2}}, \qquad R^2 = 1 - \frac{\text{MSE}}{\text{MSE}_{\text{baseline}}}\]

so \(R^2 = 0\) is no-skill and \(R^2 = 1\) is exact. Both quantities are
\textbf{column-balanced} by default: per-column errors are aggregated
with equal weight per benchmark rather than per cell, because
cell-weighting lets a handful of densely-evaluated famous benchmarks
dominate the score and renders it nearly insensitive to densification.
\(R^2\) is computed as a \emph{single pooled ratio} of column-balanced
error to column-balanced baseline, never as an average of per-column
\(R^2\) values, which is unstable when a thin column has a small
baseline. Methods that do not natively predict cells still report this
metric, by predicting each held-out cell from the row's surviving
observed cells via the conditional-Gaussian (best linear) predictor
implied by their recovered correlation matrix. Definitions and the exact
estimator for each family are in
\href{Appendix-Methods\#I\%20Held-out\%20metric}{Appendix I}.

This metric is used for hyperparameter selection within each method
(rank, \(k\), number of trees, number of imputations), and as a
\textbf{gate}: a completed matrix whose held-out \(R^2\) falls below 0.4
is not factored at all, on the grounds that a factor structure extracted
from predictions no better than a column mean is not interpretable.
Gating is reported alongside the results, so that a method's failure to
clear it is visible rather than silently absent.

\hypertarget{factor-analysis}{%
\paragraph{Factor analysis}\label{factor-analysis}}

Factoring is \textbf{held identical across every completed matrix}, so
that the imputation input is the only thing that varies. We use
minimum-residual exploratory factor analysis with promax (oblique)
rotation on the correlation matrix of the completed or surrogate data.
Rotation is oblique by design: the whole question at issue is whether
the first-order domain factors \(F\) correlate strongly enough to imply
a higher-order \(G\), and an orthogonal rotation would answer that by
assumption.

The number of factors is chosen by \textbf{Horn's parallel analysis}:
the observed eigenvalues of the correlation matrix are compared
position-by-position against the 95th percentile of eigenvalues from
random Gaussian matrices of the same shape, and the factor count is the
number of positions where the observed value exceeds the random cutoff.
Because the random baseline depends only on the matrix shape and not on
the data, it is computed once per shape and cached; it is emphatically
\emph{not} a global cutoff, and reusing one across shapes would be
invalid. The retained count is capped at 20 for tractability and further
capped at the numeric rank the matrix can actually support; where the
estimator fails at the chosen count it is decremented until it succeeds,
and the count actually used is reported
(\href{Appendix-Methods\#J\%20Factor\%20analysis\%20details}{Appendix J}).

Because the rotation is oblique the first-order factors are correlated,
and that correlation is the object of interest. We therefore decompose
the hierarchy with a \textbf{Schmid--Leiman bifactor transformation}, in
which every benchmark loads directly on the global general factor \(G\)
and on its domain factor \(F\) --- an estimator of exactly the
decomposition \(B_i = F_i + G_i + U_i\) set out in
\url{Appendix\#Definitions}. This yields three reported quantities:
\(\omega_h\), the proportion of total score variance attributable to
\(G\); \(\omega_{total}\), the proportion attributable to all common
factors; and the vector \(\omega_{hs}\) of domain-factor reliabilities.
The ratio \(\omega_h / \omega_{total}\) reads as ``of the common
variance, how much is general'' --- the direct analogue of the
\(g\)-saturation question in human psychometrics, and the quantity on
which our central claim turns. We note that this is an exploratory
Schmid--Leiman solution, not a confirmatory bifactor model with
cross-loadings constrained to zero, and we interpret it accordingly; the
contrast with Ilica \& Gignac's (2024) confirmatory approach is
deliberate.

Every cell of the design is factored \textbf{twice}: once at the
parallel-analysis factor count, and once \textbf{forced to two factors}.
The forced-2 run exists because the parallel-analysis count is itself
unstable under this data's sparsity --- a finding we report rather than
hide --- and fixing the count makes the \(G\)-related quantities
comparable across cells whose data-driven counts differ, at the cost of
certainly under-factoring some of them.

\hypertarget{implementation}{%
\paragraph{Implementation}\label{implementation}}

Corpus construction, densification, and analysis scripts are implemented
in Python; imputation and factor analysis in R, with the OneSidedMC
estimator in Julia invoked as a subprocess and its output routed through
the identical R factoring path. All results are persisted to a single
relational store keyed by (method, dataset, run), so the full design is
queryable rather than reconstructed from file names, and environments
are pinned per language.
\href{Appendix-Methods\#K\%20Software\%20environment\%20and\%20reproduction}{Appendix K}
lists exact packages and the commands that reproduce every table above.

\hypertarget{analysis}{%
\section{Analysis}\label{analysis}}

\hypertarget{robustness-and-sensitivity-analyses}{%
\paragraph{Robustness and sensitivity
analyses}\label{robustness-and-sensitivity-analyses}}

The design is a cross-product of \{3 densifiers + raw\} × \{2 collapse
strategies\} × \{completion methods\}
(\url{Methodology\#Sparsity\%20Handling}), and its cells are compared
along four axes:

\begin{enumerate}
\def\labelenumi{\arabic{enumi}.}
\tightlist
\item
  \textbf{Across densifiers and collapse strategies.} Since C, S, and R
  embody different selection biases over the same underlying corpus, and
  the two collapse strategies embody different model-identity
  assumptions, systematic variation across them is the closest available
  proxy for sensitivity to the MNAR mechanism itself.
\item
  \textbf{Across completion methods.} Where two solutions have the same
  number of factors, we compute pairwise factor congruence
  (sign-invariant cosine similarity between loading vectors, matched
  after sorting by explained variance), grouped by dataset, solution
  type, and shape, so that only like-for-like solutions are compared
  (\href{Appendix-Methods\#J.6\%20Cross-method\%20congruence}{Appendix J.6}).
\item
  \textbf{Across held-out splits.} An optional seed sweep refits each
  method across repeated random holdouts, reporting the distribution of
  held-out error, of the selected hyperparameter, and of \(\omega_h\).
  We stress that this quantifies \emph{split} variance under a fixed
  missingness pattern; it does \textbf{not} measure robustness to the
  MNAR mechanism, since every split inherits the same selection.
\item
  \textbf{Per-benchmark influence.} A leave-one-covariate-out analysis
  recomputes \(\omega_h\) with each benchmark removed in turn, giving
  \(\Delta\omega_h\) per benchmark
  (\href{Appendix-Methods\#J.5\%20Leave-one-covariate-out}{Appendix J.5}).
  Correlating \(\Delta\omega_h\) with each benchmark's observation
  frequency tests directly whether the general factor is driven by
  capability structure or by evaluation volume --- the concern raised by
  the highly unequal source composition in Table 1 of \url{Methodology},
  and the same common-investment bias we raise against prior work in the
  introduction.
\end{enumerate}

Three commitments follow and constrain how we read our own results.
First, \textbf{held-out accuracy is not evidence of factor validity}: a
method can predict held-out cells well and still deliver an unstable
rotation, so our findings rest on agreement across methods and stability
across design cells, not on predictive fit. Second, \textbf{the factor
count is weakly identified} on data this sparse, and we report that
instability as a result rather than resolving it by fiat. Third, a
factor that is dominated by a single benchmark family is a fact about
the corpus, not about language models; the influence analysis above
exists to make such cases visible.

\hypertarget{results}{%
\section{Results}\label{results}}

\hypertarget{discussion}{%
\section{Discussion}\label{discussion}}

\hypertarget{the-g-factor-is-only-partially-coherent}{%
\subsection{The g-factor is only partially
coherent}\label{the-g-factor-is-only-partially-coherent}}

\hypertarget{bottom-up-theories-of-llm-intelligence}{%
\subsection{Bottom-up theories of LLM
intelligence}\label{bottom-up-theories-of-llm-intelligence}}

Importantly what we find is that, unlike theories of human psychology ()
where intelligence are explicitly assumed to possess higher-order
structure, with fluid intelligence as a higher-order factor that have
theoretical causal effects () to the performance of other cognitive
abilities, fluid intelligence, wikipedia memorization, agentic tool use,
are more likely to be horizontal with respect to each other with no
clear higher-order structure. Ablation studies provide experimental
evidence to this point: \ldots{}

This is also not to say that LLMs are not intelligent (in our private
view, they very clearly are). But it is important to realize that just
as humans and ants are differentially but equally ``cognitive'', LLMs
are intelligent in a considerably different way to humans. This means
that a theory of intelligence requires a blank-slate, bottom-up
empirical approach without overfitting theories of human intelligence
into systems that possess significant architectural and functional
divergence from humans.

\hypertarget{mechanisms-for-a-possible-g-factor}{%
\subsection{Mechanisms for a possible
g-factor}\label{mechanisms-for-a-possible-g-factor}}

\hypertarget{references}{%
\section{References}\label{references}}

\hypertarget{psychometric-g-factor}{%
\subsection{Psychometric g-factor}\label{psychometric-g-factor}}

\begin{itemize}
\tightlist
\item
  Waterhouse, L. (2023). Why multiple intelligences theory is a
  neuromyth. \emph{Frontiers in psychology}, \emph{14}, 1217288.
\item
  Pokropek, A., Marks, G. N., \& Borgonovi, F. (2022). How much do
  students' scores in PISA reflect general intelligence and how much do
  they reflect specific abilities? \emph{Journal of Educational
  Psychology}, \emph{114}(5), 1121.
\item
  Johnson, W., Bouchard Jr, T. J., Krueger, R. F., McGue, M., \&
  Gottesman, I. I. (2004). Just one g: Consistent results from three
  test batteries. \emph{Intelligence}, \emph{32}(1), 95-107.
\item
  Johnson, W., Te Nijenhuis, J., \& Bouchard Jr, T. J. (2008). Still
  just 1 g: Consistent results from five test batteries.
  \emph{Intelligence}, \emph{36}(1), 81-95.
\item
  Van der Maas, H. L., Kan, K. J., \& Borsboom, D. (2014). Intelligence
  is what the intelligence test measures. Seriously. \emph{Journal of
  Intelligence}, \emph{2}(1), 12-15.
\item
  Major, J. T., Johnson, W., \& Bouchard Jr, T. J. (2011). The
  dependability of the general factor of intelligence: Why small,
  single-factor models do not adequately represent g.
  \emph{Intelligence}, \emph{39}(5), 418-433.
\item
  Kranzler, J. H., \& Jensen, A. R. (1989). Inspection time and
  intelligence: A meta-analysis. \emph{Intelligence}, \emph{13}(4),
  329-347.
\item
  Jensen, A. R. (2002). Psychometric g: Definition and substantiation.
  In \emph{The general factor of intelligence} (pp.~51-66). Psychology
  Press.
\end{itemize}

\hypertarget{intelligence-in-aimlcompsci}{%
\subsection{Intelligence in
AI/ML/Compsci}\label{intelligence-in-aimlcompsci}}

\begin{itemize}
\tightlist
\item
  Morris, M. R., Sohl-Dickstein, J., Fiedel, N., Warkentin, T., Dafoe,
  A., Faust, A., \ldots{} \& Legg, S. (2023). Levels of AGI for
  Operationalizing Progress on the Path to AGI. \emph{arXiv preprint
  arXiv:2311.02462}.
\item
  Chollet, F., Knoop, M., Kamradt, G., Landers, B., \& Pinkard, H.
  (2025). Arc-agi-2: A new challenge for frontier ai reasoning systems.
  \emph{arXiv preprint arXiv:2505.11831}.
\end{itemize}

\hypertarget{prior-work}{%
\subsection{Prior work}\label{prior-work}}

\begin{itemize}
\tightlist
\item
  Ilica, D., \& Gignac, G. E. (2024). Evidence of interrelated
  cognitive-like capabilities in large language models: Indications of
  artificial general intelligence or achievement? \emph{Intelligence},
  \emph{106}, 101858.
\item
  Burnell, R., Hao, H., Conway, A. R., \& Orallo, J. H. (2023).
  Revealing the structure of language model capabilities. \emph{arXiv
  preprint arXiv:2306.10062}.
\end{itemize}

\hypertarget{methods-and-statistical-tools}{%
\subsection{Methods and statistical
tools}\label{methods-and-statistical-tools}}

\begin{itemize}
\tightlist
\item
  Horn, J. L. (1965). A rationale and test for the number of common
  factors. \emph{Psychometrika}, \emph{30}(2), 179-185.
\item
  Schmid, J., \& Leiman, J. M. (1957). The development of hierarchical
  factor solutions. \emph{Psychometrika}, \emph{22}(1), 53-61.
\item
  Revelle, W., \& Zinbarg, R. E. (2009). Coefficients alpha, beta,
  omega, and the glb: Comments on Sijtsma. \emph{Psychometrika},
  \emph{74}(1), 145-154.
\item
  Revelle, W. (2024). \emph{psych: Procedures for Psychological,
  Psychometric, and Personality Research}. R package.
\item
  Lorenzo-Seva, U., \& ten Berge, J. M. (2006). Tucker's congruence
  coefficient as a meaningful index of factor similarity.
  \emph{Methodology}, \emph{2}(2), 57-64.
\item
  Mazumder, R., Hastie, T., \& Tibshirani, R. (2010). Spectral
  regularization algorithms for learning large incomplete matrices.
  \emph{Journal of Machine Learning Research}, \emph{11}, 2287-2322.
\item
  Cao, Y., Liang, Y., \& Valiant, G. (2023). One-sided matrix completion
  from two observations per row. \emph{ICML}.
\item
  Keshavan, R. H., Montanari, A., \& Oh, S. (2010). Matrix completion
  from a few entries. \emph{IEEE Transactions on Information Theory},
  \emph{56}(6), 2980-2998.
\item
  Chatterjee, S. (2015). Matrix estimation by universal singular value
  thresholding. \emph{The Annals of Statistics}, \emph{43}(1), 177-214.
\item
  Stekhoven, D. J., \& Bühlmann, P. (2012). MissForest ---
  non-parametric missing value imputation for mixed-type data.
  \emph{Bioinformatics}, \emph{28}(1), 112-118.
\item
  van Buuren, S., \& Groothuis-Oudshoorn, K. (2011). mice: Multivariate
  imputation by chained equations in R. \emph{Journal of Statistical
  Software}, \emph{45}(3), 1-67.
\item
  Higham, N. J. (2002). Computing the nearest correlation matrix --- a
  problem from finance. \emph{IMA Journal of Numerical Analysis},
  \emph{22}(3), 329-343.
\item
  Rubin, D. B. (1976). Inference and missing data. \emph{Biometrika},
  \emph{63}(3), 581-592.
\item
  Liang, P., Bommasani, R., Lee, T., et al.~(2023). Holistic evaluation
  of language models (HELM). \emph{TMLR}.
\item
  Fourrier, C., Habib, N., Lozovskaya, A., Szafer, K., \& Wolf, T.
  (2024). Open LLM Leaderboard v2. Hugging Face.
\item
  Chiang, W. L., Zheng, L., Sheng, Y., et al.~(2024). Chatbot Arena: An
  open platform for evaluating LLMs by human preference. \emph{ICML}.
\end{itemize}

%%%%%%%%%%%%%%%%%%%%%%%%%%%%%%%%%%%%%%%%%%%%%%%%%%%%%%%%%%%%

\bibliography{iclr2027_conference}
\bibliographystyle{iclr2027_conference}

\end{document}
